% =============================================================================
% Dedication — Rosalind Franklin
% =============================================================================
\clearpage
\thispagestyle{empty}
\vspace*{\fill}
\begin{center}
{\large\itshape In tribute to Rosalind Franklin}\\[0.3em]
{\small 25 July 1920 --- 16 April 1958}
\end{center}
\vspace{1.2em}

\noindent
This is a book about what counts as evidence, so it is dedicated to the person
who produced the decisive evidence for the structure of DNA and was not credited
for it.

\medskip\noindent
Rosalind Franklin was a physical chemist who had refined X-ray crystallography in
Paris before joining King's College London in 1951. In May 1952 she and her
student Raymond Gosling exposed the photograph catalogued as~51: a fibre
diffraction pattern of the B~form of deoxyribonucleic acid. It is reproduced, by
computation rather than by scan, as Figure~\ref{fig:photo51} of this book.

\medskip\noindent
What makes Photo~51 decisive is that it does not require interpretation so much
as reading. The X-shaped cross of reflections is the signature of a helix, and
the angle of the cross fixes the ratio of pitch to radius. The strong meridional
reflection at $1/3.4\,\text{\AA}^{-1}$ gives the spacing between stacked bases.
And the fourth layer line is \emph{missing} --- a systematic absence that follows
if, and essentially only if, the structure has two strands offset along the axis
by three-eighths of a turn. Franklin's notebooks record the helical parameters
and the two-chain, phosphates-outward arrangement. The data were not suggestive;
they were constraining.

\medskip\noindent
In January 1953 her colleague Maurice Wilkins showed the photograph to James
Watson without her knowledge or consent. An unpublished report of her
measurements reached Francis Crick by a similar route. The double-helix model
published in April 1953 rested on those data. Its acknowledgements did not say
so.

\medskip\noindent
Franklin died of ovarian cancer in 1958, aged thirty-seven, having moved to
Birkbeck and done foundational work on the structure of tobacco mosaic virus and
poliovirus. The Nobel Prize went to Watson, Crick and Wilkins in 1962. It is not
awarded posthumously, so the question of whether she would have shared it is
unanswerable --- but it was never put.

\medskip\noindent
The relevance to a textbook on formal methods is direct, and it is not
sentimental. This book argues that a claim should be traceable to the measurement
that supports it, that the provenance of data is part of its meaning, and that
the person who did the work should be visible in the record. Those are the same
principle stated three ways. Franklin's case is what it looks like when the
principle fails: the evidence was sound, the inference drawn from it was correct,
and the attribution was wrong --- and the attribution being wrong is itself a
defect in the scientific record, not merely an injustice to one person.

\medskip\noindent
Every figure in this book names the script that produced it and the endpoint its
data came from. That convention is the smallest possible tribute, and the only
one that is actually in our power to pay.

\medskip\noindent
There is one further debt, of method rather than of credit. What made Photo~51
decisive was partly a \emph{missing} reflection, and an absence is evidence only
when the method could have shown the thing. Section~\ref{sec:gaps} of this book
rests on the same move: its central result is a negative one, and it is worth
reporting only because the same test detects a positive control
overwhelmingly. Franklin read an absence correctly because her plate showed the
neighbouring layer lines plainly. We have tried to earn our absences the same
way.

\vspace*{\fill}
\clearpage
